% from ega4-8-fr.tex
\oldpage[IV-3]{18}
\item[\emph{b)}] Il existe $\alpha$ tel que l'espace sous-jacent à
$S_\alpha$ soit noethérien, et pour $\mu\geq\lambda$, $u_{\lambda\mu}$ est
un homéomorphisme de $S_\mu$ sur un sous-espace de $S_\lambda$.
\end{enumerate}
Sous ces conditions, il existe un $\lambda$ tel que, pour tout
$\mu\geq\lambda$~:
\begin{enumerate}
\item[(i)] Pour toute composante irréductible $Y_i$ de $S$
$(1\leq i\leq m)$, $\overline{u_\mu(Y_i)}$ est une composante irréductible
de $S_\mu$, et l'application $Y_i\mapsto\overline{u_\mu(Y_i)}$ est une
bijection de l'ensemble des composantes irréductibles de $S$ sur l'ensemble
des composantes irréductibles de $S_\mu$.
\item[(ii)] Pour toute composante connexe $C_j$ de $S$
$(1\leq j\leq n)$, $\overline{u_\mu(C_j)}$ est une composante connexe de
$S_\mu$, et l'application $C_j\mapsto\overline{u_\mu(C_j)}$ est une
bijection de l'ensemble des composantes connexes de $S$ sur l'ensemble des
composantes connexes de $S_\mu$.
\end{enumerate}
\end{proposition}

Nous établirons d'abord le

\begin{lemma}[8.4.2.1]
\label{IV.8.4.2.1-fr}
Sous la condition \emph{a)} ou \emph{b)} de (8.4.2), il existe $\lambda$ tel
que, pour $\mu\geq\lambda$, $u_\mu:S\to S_\mu$ soit dominant.
\end{lemma}

Dans le cas \emph{a)}, cela a déjà été prouvé sans supposer l'espace $S$
noethérien (8.3.8, (i)). Dans le cas \emph{b)}, posons
$Z_\alpha=\overline{u_\alpha(S)}$~; en tant que partie fermée de l'espace
noethérien $S_\alpha$, $Z_\alpha$ est constructible, et comme
$u_\alpha^{-1}(Z_\alpha)=S$, il résulte de (8.3.5) qu'il existe
$\lambda\geq\alpha$ tel que, pour $\mu\geq\lambda$, on ait
$Z_\mu=u_{\alpha\mu}^{-1}(Z_\alpha)=S_\mu$. Mais comme $u_{\alpha\mu}$ est
un homéomorphisme de $S_\mu$ sur un sous-espace de $S_\alpha$, et comme
l'application composée $S\to Z_\mu\to Z_\alpha$ est dominante, il en est de
même de $S\to Z_\mu=S_\mu$.

Ce lemme étant démontré, on peut supposer que pour tout $\lambda$,
$u_\lambda$ est un morphisme dominant.

\medskip
\noindent
(i) Chacun des $S_\lambda$ est réunion des
$\overline{u_\lambda(Y_i)}$, qui sont irréductibles. D'autre part, si $U_i$
est l'ouvert de $S$ complémentaire de la réunion des $Y_j$ d'indice $j\neq i$
$(1\leq i\leq m)$, les $U_i$ sont deux à deux disjoints et
$Y_i=\overline{U_i}$ $(0_{\mathrm I},2.1.6)$. Comme l'espace sous-jacent à
$S$ est noethérien, les $U_i$ sont quasi-compacts, donc il existe un indice
$\lambda$ et des ouverts $U_{i\lambda}$ de $S_\lambda$ tels que
$U_i=u_\lambda^{-1}(U_{i\lambda})$ pour $1\leq i\leq m$ (8.2.11). On en
conclut que si l'on pose
$U_{i\mu}=u_{\lambda\mu}^{-1}(U_{i\lambda})$ pour $\lambda\leq\mu$, les
$U_{i\mu}$ sont deux à deux disjoints, car les
$U_i=u_\mu^{-1}(U_{i\mu})$ le sont et $u_\mu$ est dominant. Par suite,
aucune des adhérences $\overline{U_{i\mu}}$ n'est contenue dans une autre et
$u_\mu(U_i)$ est dense dans $U_{i\mu}$ puisque $u_\mu$ est dominant~; on a
donc
$\overline{U_{i\mu}}=\overline{u_\mu(Y_i)}$, ce qui prouve que les
$\overline{U_{i\mu}}$ sont les composantes irréductibles de $S_\mu$
$(0_{\mathrm I},2.1.7)$ et achève la démonstration.

\medskip
\noindent
(ii) Comme l'espace $S$ est noethérien, les $C_j$ sont ouverts et fermés dans
$S$ et quasi-compacts~; le même raisonnement que dans (i) montre donc qu'il
existe un $\lambda$ et des ouverts $V_{j\lambda}$ de $S_\lambda$ tels que
$C_j=u_\lambda^{-1}(V_{j\lambda})$ pour $1\leq j\leq n$. On voit aussi,
comme dans (i), que si l'on pose
$V_{j\mu}=u_{\lambda\mu}^{-1}(V_{j\lambda})$ pour $\mu\geq\lambda$, les
$V_{j\mu}$ sont deux à deux disjoints, et $u_\mu(C_j)$ est dense dans
$V_{j\mu}$~; cela entraîne que $V_{j\mu}$ est connexe. En outre, il résulte
de (8.3.4) que pour $\mu$ assez grand, la réunion des $V_{j\mu}$ est
$S_\mu$, puisque tout ouvert dans un préschéma est ind-constructible (1.9.6).
Les $V_{j\mu}$ sont donc les composantes connexes de $S_\mu$, ce qui achève
la démonstration.

On notera que si les morphismes $u_{\lambda\mu}$ sont des \emph{immersions},
ils vérifieront en particulier la condition \emph{b)} de (8.4.2).

